\documentclass[../main.tex]{subfiles}

\begin{document}

\chapter{V.2 信息、结构与正交性}
\label{math:structure}

\begin{motivation}
	\textbf{你有没有想过}，为什么我们的思考中总是充满冗余？如何用一个数学过程系统地"去重"，提取出思想的真正正交基？

	\textbf{核心动机}：将 Gram--Schmidt 正交化过程作为逻辑去重的隐喻，展示如何从冗余的思维向量中提取最小正交基，从而理解"结构"的本质就是去重。
\end{motivation}

\begin{logicmap}
\begin{tikzpicture}[node distance=1.2cm, every node/.style={draw, rectangle, rounded corners, align=center}]
\node (problem) {冗余思维：多个向量不正交};
\node (insight) [below left=of problem] {核心洞见：思想=向量，废话=投影};
\node (method) [below right=of problem] {Gram--Schmidt 正交化};
\node (result) [below=of method] {正交基：去重后的最小结构};
\draw[-{Stealth}] (problem) -- (insight);
\draw[-{Stealth}] (problem) -- (method);
\draw[-{Stealth}] (insight) -- (result);
\draw[-{Stealth}] (method) -- (result);
\end{tikzpicture}
\end{logicmap}

\mathlink{math:structure}{正文链接：结构为何等价于去重？}

\section{逻辑去重的 Gram--Schmidt 过程}

你有没有想过，为什么我们的思考中总是充满冗余？为什么同一个问题，我们用不同的方式反复表达，却很少意识到这些表达本质上是相关的？\\
如果思维是一组向量，我们如何用数学方法系统地"去重"，提取出真正独立的、正交的核心概念？

\begin{mathbox}{问题动机与写作映射}
	思想 = 向量，废话 = 投影
\end{mathbox}

\begin{mathbox}{数学过程}
	\[
	u_k = v_k - \sum_{j<k} \text{proj}_{u_j}(v_k)
	\]
\end{mathbox}

\subsection{逻辑衔接}
通过 Gram--Schmidt 正交化，我们展示了如何从冗余的思维向量中提取最小正交基，这正是"结构就是去重"的数学体现。这种将复杂系统化简为独立维度的思想，将在下一章「历史的条件概率」中继续深化：当我们面对历史的多因素影响时，正交基提供了识别真正独立变量的工具，帮助我们在不确定性和相关性中找到可分析的结构。
\end{document}
